%--------------------------------------
% Caso de Uso: Eliminar Servicio de Reservación
%--------------------------------------

\begin{UseCase}{CU19}{Eliminar Servicio de Reservación}{
	Permitir que un usuario autorizado elimine o cancele un servicio adicional previamente agregado a una reservación.
}
	\UCitem{Versión}{\color{Gray}1.0}
	\UCitem{Autor}{\color{Gray}Equipo de Desarrollo}
	\UCitem{Supervisa}{\color{Gray}Coordinador del Sistema}
	\UCitem{Actor}{\hyperlink{Administrador}{Administrador} o \hyperlink{Empleado}{Empleado}}
	\UCitem{Propósito}{Remover un servicio adicional de una reservación, actualizando costos y disponibilidades.}
	\UCitem{Entradas}{Servicio seleccionado y motivo opcional.}
	\UCitem{Origen}{Pantalla de detalle de reservación}
	\UCitem{Salidas}{Servicio cancelado, costo actualizado y disponibilidad liberada.}
	\UCitem{Destino}{Pantalla de detalle de reservación y base de datos}
	\UCitem{Precondiciones}{
		\begin{itemize}
			\item El usuario debe haber iniciado sesión con permisos adecuados.
			\item La reservación debe existir.
			\item El servicio debe estar asociado a la reservación.
			\item El servicio debe encontrarse en un estado cancelable.
		\end{itemize}
	}
	\UCitem{Postcondiciones}{
		\begin{itemize}
			\item El servicio queda cancelado o eliminado.
			\item El costo total de la reservación se actualiza.
			\item El empleado asignado queda disponible.
			\item El horario del servicio se libera en el calendario.
		\end{itemize}
	}
	\UCitem{Errores}{
		Servicio no encontrado, permisos insuficientes, servicio no cancelable, error al actualizar la información.
	}
	\UCitem{Tipo}{Caso de uso secundario}
	\UCitem{Observaciones}{
		\begin{itemize}
			\item Los servicios ya realizados no pueden eliminarse.
			\item Puede aplicarse un cargo por cancelación si el servicio estaba próximo a realizarse.
			\item Se notifica al empleado asignado cuando el servicio es cancelado.
		\end{itemize}
	}
\end{UseCase}

%--------------------------------------
\begin{UCtrayectoria}
	\UCpaso[\UCactor] Accede a la pantalla de detalle de una reservación.
	\UCpaso[\UCactor] Selecciona un servicio asociado a la reservación.
	\UCpaso[\UCactor] Selecciona la opción para eliminar el servicio.
	\UCpaso El sistema solicita confirmación y permite ingresar un motivo opcional.
	\UCpaso[\UCactor] Confirma la eliminación del servicio.
	\UCpaso El sistema verifica los permisos del usuario.
	\UCpaso El sistema valida que el servicio pueda cancelarse.
	\UCpaso El sistema registra la cancelación del servicio.
	\UCpaso El sistema libera la disponibilidad del empleado asignado.
	\UCpaso El sistema actualiza el costo total de la reservación.
	\UCpaso El sistema libera el horario en el calendario de servicios.
	\UCpaso El sistema muestra un mensaje indicando que el servicio fue eliminado correctamente.
	\UCpaso El sistema actualiza la pantalla de detalle de la reservación.
\end{UCtrayectoria}

%--------------------------------------
\begin{UCtrayectoriaA}{A}{Servicio no encontrado}
	\UCpaso El sistema informa que el servicio seleccionado no existe.
	\UCpaso El caso de uso termina sin realizar cambios.
\end{UCtrayectoriaA}

%--------------------------------------
\begin{UCtrayectoriaA}{B}{Permisos insuficientes}
	\UCpaso El sistema informa que el usuario no tiene permisos para eliminar el servicio.
	\UCpaso El caso de uso termina sin realizar cambios.
\end{UCtrayectoriaA}

%--------------------------------------
\begin{UCtrayectoriaA}{C}{Servicio no cancelable}
	\UCpaso El sistema informa que el servicio no puede eliminarse en su estado actual.
	\UCpaso El caso de uso termina sin realizar cambios.
\end{UCtrayectoriaA}

%--------------------------------------
\subsection{Puntos de extensión}
\UCExtenssionPoint{
	El usuario desea ajustar cargos derivados de la cancelación del servicio.
}{
	Del paso 8.
}{
	Procesar Ajuste por Cancelación.
}

%--------------------------------------
% Fin del Caso de Uso
